\documentclass[../../../uem_phy_std_q.tex]{subfiles}

\begin{document}
    図のように，真空中において，同形の4枚の金属極板X，Y，Z，Wがあり，
    XとYが間隔$2d$で，ZとWが間隔$d$で平行に向き合わせてある．
    それらは，電池B（起電力$V_0$），スイッチ$\mathrm{S}_1$，$\mathrm{S}_2$，
    およびアースと接続されて回路を形成している．
    XとYからなる平行平板コンデンサの静電容量を$C_0$とし，端の効果は無視する．
    はじめ，$\mathrm{S}_1$，$\mathrm{S}_2$は共に開いており，いずれの極板も帯電していない．

    \begin{enumerate}
        \item $\mathrm{S}_1$を閉じて安定するまで待った．%
            このとき，X下面の電荷$x_1$，Z下面の電荷$z_1$を求めよ．%
            また，XY間の電場$E_1$（X→Yの向き正）を求めよ．
        \item $\mathrm{S}_1$を開き，$\mathrm{S}_2$を閉じて安定するまで待った．%
            このとき，X下面の電荷$x_2$，Z下面の電荷$z_2$を求めよ．%
            また，XY間の電場$E_2$（X→Yの向き正），Xの電位$\phi_2$を求めよ．
        \item $\mathrm{S}_2$を開き，$\mathrm{S}_1$を閉じて安定するまで待った．%
            このとき，X下面の電荷$x_3$，Z下面の電荷$z_3$を求めよ．
        \item $\mathrm{S}_1$を閉じたまま，$\mathrm{S}_2$も閉じて安定するまで待った．%
            このとき，X下面の電荷$x_4$，Z下面の電荷$z_4$を求めよ．
        \item $\mathrm{S}_2$を閉じたまま，$\mathrm{S}_1$を開き，%
            ZとWの間を誘電体（比誘電率$\varepsilon_r$）で満たし，安定するまで待った．%
            このとき，X下面の電荷$x_5$，Z下面の電荷$z_5$を求めよ．
    \end{enumerate}
    \begin{minipage}{\linewidth}
        \centering
            \includegraphics{\subfix{fig_uem_phy_std_em_cb_12/fig_uem_phy_std_em_cb_12_00_q.pdf}}
    \end{minipage}
\end{document}




